\documentclass[11pt]{article}
\usepackage{amsmath,amssymb,amsthm}
\usepackage[margin=1in]{geometry}

\newtheorem{theorem}{Theorem}
\newtheorem{lemma}[theorem]{Lemma}

\title{Problem 515: Dissonant Numbers}
\author{Project Euler Solutions}
\date{}

\begin{document}
\maketitle

\section{Problem Statement}

For a positive integer $n$ and base $b \geq 2$, let $\mathrm{ds}_b(n)$ denote the digit sum of $n$ in base~$b$. Define $D(b, s, N)$ as the count of integers $1 \leq n \leq N$ with $\mathrm{ds}_b(n) = s$. Compute $\sum D(b_i, s_i, N_i)$ for given parameter triples.

\section{Mathematical Foundation}

\begin{theorem}[Generating Function for Digit Sums]
The number of non-negative integers with exactly $L$ base-$b$ digits (allowing leading zeros) whose digit sum equals $s$ is
\[
[x^s]\left(\frac{1 - x^b}{1 - x}\right)^L.
\]
\end{theorem}

\begin{proof}
Each digit $d_i \in \{0, \ldots, b-1\}$ has generating function $\sum_{d=0}^{b-1}x^d = \frac{1 - x^b}{1 - x}$. Independence of the $L$ digits gives the $L$-fold product. The coefficient of $x^s$ counts strings with digit sum~$s$.
\end{proof}

\begin{lemma}[Inclusion-Exclusion]
\[
[x^s]\left(\frac{1-x^b}{1-x}\right)^L = \sum_{j=0}^{\lfloor s/b \rfloor}(-1)^j \binom{L}{j}\binom{s - jb + L - 1}{L - 1}.
\]
\end{lemma}

\begin{proof}
Expand $(1-x^b)^L = \sum_j (-1)^j \binom{L}{j}x^{jb}$ and $(1-x)^{-L} = \sum_m \binom{m+L-1}{L-1}x^m$. The coefficient of $x^s$ in their product is obtained by convolution, yielding the stated sum.
\end{proof}

\begin{theorem}[Digit DP Correctness]
The digit DP with states $(\mathrm{pos}, \mathrm{rem}, \mathrm{tight})$ correctly counts all integers $n \in [0, N]$ with $\mathrm{ds}_b(n) = s$.
\end{theorem}

\begin{proof}
Write $N = (d_k \cdots d_0)_b$. The DP constructs $n$ digit-by-digit from position~$k$ down to~$0$. The \texttt{tight} flag ensures $n \leq N$: when tight, the current digit is bounded by $d_{\mathrm{pos}}$; choosing any smaller digit releases the constraint for all subsequent positions. The \texttt{rem} parameter tracks the remaining digit sum. Each integer $n \leq N$ corresponds to exactly one DP path, so the count is exact.
\end{proof}

\section{Algorithm}

\begin{enumerate}
    \item Write $N$ in base $b$: $N = (d_k \, d_{k-1} \, \cdots \, d_0)_b$.
    \item Initialize DP: $\mathrm{dp}(\mathrm{pos}, \mathrm{rem}, \mathrm{tight})$.
    \item At each position, choose digit $d \in \{0, \ldots, \min(b-1, d_{\mathrm{pos}} \cdot [\mathrm{tight}])\}$.
    \item Transition: $\mathrm{dp}(\mathrm{pos}+1, \mathrm{rem}-d, \mathrm{tight} \land (d = d_{\mathrm{pos}}))$.
    \item Base case: at $\mathrm{pos} = L$, return $[\mathrm{rem} = 0]$.
    \item Subtract the $n = 0$ case if $s = 0$.
\end{enumerate}

\section{Complexity Analysis}

\begin{itemize}
    \item \textbf{Time:} $O(L \cdot s \cdot b)$ per query, where $L = \lfloor \log_b N \rfloor + 1$.
    \item \textbf{Space:} $O(L \cdot s)$ for memoization.
\end{itemize}

\section{Answer}

\[
\boxed{2422639000800}
\]

\end{document}
